\noindent \textbf{Lenke til all kode og alle plotter for innleveringen finnes i Vedlegg B.}

\hypertarget{oppgave1}{}
\section*{Oppgave 1}
\addcontentsline{toc}{section}{Oppgave 1}

% --------------------------------------------------
\subsection*{a)}
Det opprettes en radvektor med partall fra 2 til 20. 
Kolonoperatoren brukes med en steglengde på 2.
\begin{verbatim}
lv_a = 2:2:20;
\end{verbatim}
Dette gir radvektoren
\[
lv_a =
\begin{bmatrix}
2 & 4 & 6 & 8 & 10 & 12 & 14 & 16 & 18 & 20
\end{bmatrix}
\]
% --------------------------------------------------

\subsection*{b)}
Det opprettes en radvektor hvor tallene fra 1 til 10 opphøyes
i tredje potens ved hjelp av \texttt{power()}:
\begin{verbatim}
lv_b = power(1:10, 3); # eller lv_b = (1:10).^3;
\end{verbatim}
Dette gir radvektoren:
\[
lv_b =
\begin{bmatrix}
1 & 8 & 27 & 64 & 125 & 216 & 343 & 512 & 729 & 1000
\end{bmatrix}
\]

% --------------------------------------------------
\subsection*{c)}
Det opprettes en vektor hvor tallene fra 1 til 10 opphøyes
i andre potens. Apostrofen på slutten transponerer resultatet
fra en radvektor til en kolonnevektor:
\begin{verbatim}
kv_c = power((1:10), 2)';
\end{verbatim}
Dette gir kolonnevektoren:
\[
kv_c =
\begin{bmatrix}
1 \\
4 \\
9 \\
16 \\
25 \\
36 \\
49 \\
64 \\
81 \\
100
\end{bmatrix}
\]
% --------------------------------------------------
\subsection*{d)}
Her hentes annenhver verdi fra \texttt{kv\_c}, med start på
element nummer 2. Resultatet transponeres deretter til en radvektor:
\begin{verbatim}
lv_d = kv_c(2:2:end)';
\end{verbatim}
Dette gir radvektoren: 
\[
    lv_d = \begin{bmatrix}
            4 & 16 & 36 & 64 & 100
        \end{bmatrix}
\]